\cotefile{78}
\scanpage{47}
trouver une notion convenable \inserted{de ``\unsure{répère|repère}''} ``d'\uline{épinglage}''
pour les $\mathrm{SR}n\mathrm{S}$, permettant de \deleted{les} ramener
à la question d'une structure épinglée, qu'on
pourrait soumettre à un calcul\ldots{} La théorie
de descente devrait alors permettre, \deleted{en}
\deleted{\unsure{depuis|effet|définissant}} comme le ``module'' a disparu,
de \deleted{descendre} trouver un objet épinglé
canonique sur le corps premier ($\mathbb{Q}$ ou
un $\mathbb{F}_p$), dont les autres se déduisent
par ``torsion'' par un \deleted{$(\unsure{G_k|G-k})$ to} $G_k$-torseur.
C'est sans doute ainsi dans le cas \inserted{exceptionnel} $G=\mathrm{Aff}(S)\simeq
(\mathbb{Z}/n\mathbb{Z})^*.(\mathbb{Z}/n\mathbb{Z})$, mais douteux dans le
cas $G=\mathrm{Aff}^{\pm}(S)$. Essentiellement, la question
revient à celle-ci : les $\mathrm{SR}n\mathrm{S}$ \deleted{qui} forment
une gerbe sur le corps premier $\mathbb{F}_p$
(si \deleted{\unsure{en}} $p\nmid n$, ou si $n=p\ne\unsure{2|p}$, \deleted{\unsure{mais si|premier}} ou si
$n=2p$), \deleted{\unsure{ou bien est-ce}} \inserted{celle-ci a-t-elle} une \uline{section} i.e.
existe-t-il un $\mathrm{SR}n\mathrm{S}$ sur le corps
premier $\mathbb{F}_p$ ? (Ici, bien sûr, on s'intéresse
aux structures non néc\textsuperscript{t}. scindées, où
\deleted{\unsure{les}} $S$ est un sous-schéma du fibré affine
associé à $E$, lequel sous-schéma peut
être sans points). Le mieux qu'on
% archive sheet number: 46
\scanpage{48}
puisse espérer dans cette direction est ceci :

a) Si $n$ \inserted{$\geq 3$} n'est pas premier $\ne 2$ ni de la forme $2p$,
$p$ premier, trouver un $\mathrm{SR}n\mathrm S$ sur $A=\mathbb Z[\frac1n]$

b) Si $n$ est premier impair ou $n=4$,
trouver un $\mathrm{SR}n\mathrm S$ sur $A=\mathbb Z$

c) Si $n=2p$, $p$ premier impair, trouver
un $\mathrm{SR}n\mathrm S$ sur $A=\mathbb Z[\frac12]$.

La solution à a) est évidente, en prenant
\[
\left\{\begin{aligned}
Z_0&=\unsure{\mu_{n\,A}|\mu_{n_A}|\mu_{n,A}}\subset\mathrm{Sl}(2)_{\mathbb Z}\subset\underline{\mathrm{Aff}}\,\mathbb E^2_{/A}\\
\sigma_0&=\begin{pmatrix}0&1\\1&0\end{pmatrix}\in\mathrm{Sl}(2,\mathbb Z)\subset\mathrm{Aff}(\mathbb E_A^2)\\
s_0&=(1,1)\in\mathbb E^2(A)^\sigma\qquad S_0=Z.s_0
\end{aligned}\right.
\]
On trouve le \deleted{\unsure{grou|genre|groupe}} $\mathrm{SR}n\mathrm S$ type $(S_0,E_0)$
dont le schéma des autom est $G=$
$\{\pm1\}_A.\mu_{n\,A}$ (produit $\frac12$ direct) -- la
cat. fibrée des $\mathrm{SR}n\mathrm S$ s'identifie alors
à celle des $G$-torseurs. Ceci
marche encore pour \uline{tt} $n\geq3$, pourvu
qu'on se place sur $\mathbb Z[1/n]$ --
donc dans les cas b) et c), on \uline{perd}
la car. $p$, où les choses sont spéciales.
Bien sûr, \deleted{sur les \unsure{co|corps}} \inserted{dans ces cas-là}, sur le corps
% archive sheet number: 47
\scanpage{49}
premier $\mathbb F_p$, il y a une \deleted{\unsure{représent|représent.|représentation}} \inserted{section} de notre
gerbe -- et il y a = un $\mathrm{SR}n\mathrm S$
scindé sur $\mathbb F_p$ ! Mais sur $\mathbb Z$ (cas b)
resp $\mathbb Z[\frac12]$ (cas c)) ?
% The stepped outline and the crosshatching cancel the following passage as a block.
\deleted{Le plus}
\deleted{simple semblerait \unsure{«} de regarder le}
\deleted{cas (a) au dessus de $\unsure{\mathrm{Sp}|\mathrm{Spec}}(A[\frac1n])=T\smallsetminus\{\unsure{p|(p)}\}=T'$,}
\deleted{(où $T=\mathrm{Spec}(A)$) avec $\unsure{\mu_{n_{T'}}|\mu_{n\,T'}|\mu_{n,T'}}\subset\mathrm{Sl}(2)_{T'}$, et}
\deleted{d'en prendre l'\unsure{adhé|adhé'|orbite}}

\deleted{\uline{Construction}. $T=\mathrm{Spec}(A)$, $\unsure{T'=T\smallsetminus p|T=T'}$}
\[
\deleted{E=\mathbb E^2_T,\quad s_0}
\]
% archive sheet number: 48
\scanpage{50}
(\unsure{36|56} p)

\begin{flushright}
Cours\\
Polygones\\
réguliers\\
$\overline{\underline{\mathrm{II}}}$
\end{flushright}
% archive sheet number: 49
\scanpage{51}
\uline{1} Automorphismes et repères des polygônes combinatoires. Polygône
combinatoire-type
\[
\begin{aligned}
\unsure{D_n|\mathbb D_n}&=\{\sigma,u\mid\sigma^2=\unsure{1\mathrm d|1},\sigma u\sigma^{-1}=u^{-1},u^n=1\}=\{\sigma,\sigma'\mid\sigma^2=\sigma'^{2}=1,(\sigma'\sigma)^n=1\}\\
D_\infty&=\{\sigma,u\mid\sigma^2=1,\sigma u\sigma^{-1}=u^{-1}\}=\{\sigma,\sigma'\mid\sigma^2=\sigma'^{2}=1\}=\mathbb Z_2*\mathbb Z_2
\end{aligned}
\]
2. Espaces affines -- Groupe affine $\mathrm{Aff}(E)$ ($\simeq\mathrm{Aff}(n,k)\simeq
\mathrm{Gl}(n).k^n$)

3. Métrique sur espaces affines et formes quadratiques.
Formes quadratiques non dégénérées.

Rappel des points développés la fois précédente :

a) Polygône géométrique comme couple $(S,A)$
\[
\begin{aligned}
S&\subset E\quad\text{(plan affine)}\\
A&\subset\mathrm{Dr}(E)
\end{aligned}
\]
Axiome : pour la relation d'incidence habituelle
$(s\in a)$, $(S,A)$ est un polygône combinatoire.

b) Réalisation géométrique d'un polygône
combinatoire $\Pi=(S,A,R)$
\diagram{\begin{tabular}{@{}c|l@{}}
\begin{tikzcd}[ampersand replacement=\&,row sep=small,column sep=large]
S\arrow[r,"\varphi_S"]\& E\\
A\arrow[r,"\varphi_A"]\& \mathrm{Dr}(E)
\end{tikzcd}&\shortstack[l]{compatibles avec\\\unsure{relation|relations} d'incidence}
\end{tabular}}
Si $\varphi_S(\unsure{s'|s})\ne\varphi_S(s)$ quand $s,s'\in S$ sont adjacents, \uline{alors}
$\varphi_A$ est connu quand on connaît $\varphi_S$ qui peut être
pris arbitrairement par ailleurs. On dira que
$\varphi$ est \deleted{\uline{\unsure{non dégénéré|non dégénérée}}} \inserted{\uline{fidèle}} si
\[
\varphi_S(s)\text{ incident à }\varphi_A(a)\Longrightarrow s\text{ incident à }a
\]
Cela implique $\varphi_S,\varphi_A$ injectifs et $\overline S=\varphi_S(S)$, $\overline A=\varphi_A(A)$
forment un \deleted{\unsure{un|et|vrai}} polygône dans $E$. On \deleted{\unsure{prétend|prétien}} \inserted{dira}
qu'on a un polygône $\Pi$-\uline{épinglé}, ou $\Pi$-\uline{polygône} (le plus souvent, $\Pi$ sera fixé, $E$ et la
% archive sheet number: 50
\scanpage{52}
\handwritten{%
réalisation dans $E$ sont variables\ldots)

b) Polygône régulier. On avait proposé \inserted{\unsure{la|de}} (dans contexte
métrique, \inserted{sur $\mathbb R$}) : les \inserted{groupe $Z$ des} \unsure{déplacement|déplacements} qui invarient le
polygône doivent être transitifs sur l'ens des
sommets. \deleted{Dans} \inserted{Nous restreignant pour simplifier} au cas fini \unsure{Ceci|Ceci il|il} est clair :
\unsure{primo|priori} qu'il y a un pt invariant -- centre de
gravité des l'ens des sommets -- \deleted{le groupe des} \inserted{les \uline{déplacements}}
\inserted{sont des} rotations, \deleted{et le groupe} le groupe \inserted{$R$} des rotations
de centre $o$ étant le ``groupe des angles'', can\textsuperscript{t}
isom. à $\mathbb R/\mathbb Z\simeq U=\{z\in\mathbb C\mid z\bar z\unsure{)|\rvert}=1\}$. On
trouvait, si $\zeta\in\unsure{R|\mathbb R}$ \deleted{est} tel que
\[
\zeta s_0=s_1
\]
\[
\text{que }\left\{\begin{aligned}
&\zeta^n=1\qquad\zeta^d\ne1\quad\text{si }1\leq d<n,\ \text{et les}\\
&S=\{\underbrace{\zeta^i s_0}_{s_i}\mid i\in\mathbb Z/n\mathbb Z\}\\
&a_i=\mathrm{Dr}(s_i,s_{i+1})
\end{aligned}\right.
\]
% slanted
\margin{en fait, $S$ et $A$ sont des \uline{torseurs} sous $Z$ i.e. $Z$ est simplement et transitif sur $S$ et $A$}
et inversement, il est clair que les polygônes
ainsi obtenus sont ``réguliers'' en le sens proposé.

Si on prend les autom \uline{euclidiens} de $E$
sans plus \inserted{qui \unsure{conservent|stabilisent|laissent} $P$}, ils ont une propriété \deleted{plus} de transitivité
plus forte : \deleted{ils sont} \inserted{leur groupe $D$ est} \unsure{transitifs|transitif} sur l'ens
des \uline{repères} du polygône (et évidemment, = simplement
transitif\unsure{\}|)}, puisque $S$ engendre affinement $E$,
}
% letterhead
% Printed upside down relative to the notes.
\begin{tabular}{@{}l@{\qquad}l@{}}
UNIVERSITE DES SCIENCES ET TECHNIQUES &\\
DU LANGUEDOC &\\
\rule{1.4em}{.35pt}& Montpellier, le\\
INSTITUT DE MATHEMATIQUES &\\
Place Eugène Bataillon &\\
34060 MONTPELLIER CEDEX &\\
\rule{1.4em}{.35pt}&
\end{tabular}
% archive sheet number: 51
\scanpage{53}
Nous préférons ne pas nous donner de métrique
d'avance, et prendre les autom. affines de $E$
qui conservent $E$, et exiger qu'ils \unsure{sont|soit|soient} transitifs
sur l'ens des repères. On \unsure{récupèrera|récupérera|récupère} la métrique
après-coup -- de façon unique modulo homothétie
-- mais parfois elle sera dégénérée ! (Si on exigeait
dès le début une métrique non dégénérée, on perdrait
des cas \deleted{\unsure{d'un grand intérêt}} \inserted{fort intéressants} du point de
vue arithmétique -- en fait, les seuls cas vraiment
nouveaux !) Pratiquement, ce nouveau pt
de vue consiste à regarder \uline{tt} \uline{triangle} comme
polygône régulier, ainsi que \uline{tout} \uline{parallélogramme} !
(Mais on \unsure{réc.|vis.|récup.|vrais} régulier \uline{pour} une métrique
\uline{donnée}\ldots)

\uline{Définition en forme}

1) Polygône régulier (\uline{NB} pour tt polygône, l'ens des
sommets engendre affinement $E$, donc un autom
affin est connu quand on connaît sa restriction à $S$)

2) Réalisation géométrique régulière d'un polygône
combinatoire $\Pi$\ldots
% archive sheet number: 52
\scanpage{54}
% letterhead
% Printed upside down relative to the archive sheet number.
\begin{tabular}{@{}l@{\qquad}l@{}}
UNIVERSITE DES SCIENCES ET TECHNIQUES &\\
DU LANGUEDOC &\\
\rule{1.4em}{.35pt}& Montpellier, le\\
INSTITUT DE MATHEMATIQUES &\\
Place Eugène Bataillon &\\
34060 MONTPELLIER CEDEX &\\
\rule{1.4em}{.35pt}&
\end{tabular}
% archive sheet number: 53
\scanpage{55}
\uline{Polygône combinatoire et repères}

\textcircled{\small 1.} Polygône combinatoire $(S,A,R)$,
\diagram{\begin{tikzcd}[ampersand replacement=\&,row sep=tiny,column sep=0pt]
{}\&\text{sommets}\&{}\&\text{arêtes}\&{}\&\text{repères}\\
S\arrow[ur]{},\&{}\& A\arrow[ur]{}\,\text{ens,}\&{}\& R\arrow[ur]{}\subset S\times A
\end{tikzcd}}
relation, axiomes \quad 1) $R\to A$ de degré 2 \quad 2) $R\to S$ de
degré 2, \quad 3) Connexité. [Quand on laisse tomber
l'axiome 3) de connexité, on parle de ``contour
combinatoire'' ; quand on laisse tomber \deleted{\unsure{$R'$|$R$|$R\subset$}} l'axiome 2)
($R\to S$ de degré 2) on parle de graphe ; quand
on \deleted{prend} \inserted{se donne} $R\to S$ et $R\to A$, sans supposer $R\hookrightarrow S\times A$,
on parle de \deleted{\unsure{graphe|graphes}} complexe \deleted{combinatoire} \inserted{1-cellulaire}
combinatoire (ou parfois encore de ``graphe'' en
un sens légèrement généralisé).

\textcircled{\small 2.} Si
\diagram{\begin{tikzcd}[ampersand replacement=\&,row sep=0pt,column sep=large]
\raisebox{-0.9ex}{$R$}\arrow[r]\& S\\
\phantom{R}\arrow[r]\& A
\end{tikzcd}}
est un contour, on définit
\[
\begin{aligned}
\sigma_1&\quad(\text{involution sans pt fixe de }R\text{ tel que }R/\sigma_1\xrightarrow{\sim}S)\\
\sigma_0&\quad(\qquad-\qquad-\qquad\qquad-\qquad-\quad R/\sigma_0\simeq A)
\end{aligned}
\]
Soit $\unsure{G_1|\mathbb G_1|\mathfrak G_1}$ le groupe défini par
\[
G_1=\{\sigma_0,\sigma_1\mid\sigma_0^2=\sigma_1^2=1\}\xrightarrow{\sim}\mathbb Z/2\mathbb Z*\mathbb Z/2\mathbb Z
\]
alors on a un foncteur allant de la
catégorie des contours dans la cat. des $G_1$-ens.
$R$ tels que \quad a) $\sigma_0,\sigma_1$ opèrent ss pt fixe dans $E$
b) $\forall x\in E$, on a $\sigma_0x\ne\sigma_1x$ [ou encore les $G_1$-orbites
dans $E$ sont de cardinal $\unsure{>|\geq|\to}2$] -- qui exprime
que $R\to S\times A$ (où $S\overset{\mathrm{dfn}}=R/\sigma_1$, $A\overset{\mathrm{dfn}}=R/\sigma_0$) est
% archive sheet number: 54
\scanpage{56}
injectif\unsure{)|]} \deleted{\unsure{-- ou encore que}}.

\uline{Prop} Ce foncteur
\diagram{\begin{tikzcd}[ampersand replacement=\&,column sep=large]
(\text{Contours \unsure{comb.|comb}})\arrow[r]\& G_1\text{-ens }\text{``admissibles''}
\end{tikzcd}}
est une équivalence des catégories.
C'est une tautologie\ldots.

\margin{\uline{NB} la condition $\sigma_0x\ne\sigma_1x\quad\forall x\in R$ signifie que l'on exclut les ``monogones'' = polygones à 1 côté\ldots}
\textcircled{\small 3} \deleted{Pour le voir, \unsure{on}} \inserted{Il est utile d'}écrire $G_1$ \inserted{$=D_\infty$} à l'aide des
générateurs \deleted{$\unsure{\sigma_0,\sigma_1|\sigma,\sigma_1|\sigma_1}$ et}
\[
\left\{\begin{aligned}\sigma&=\sigma_1\\u&=\deleted{\unsure{\sigma_2|\sigma_1}}\sigma_1\sigma_0\end{aligned}\right.
\quad\Longrightarrow\quad
\begin{aligned}\sigma_1&=\sigma\\\sigma_0&=\deleted{\unsure{u|1|\sigma}}\sigma u=u^{-1}\sigma\end{aligned}
\]
alors on a
\[
G_1=\{\sigma,u\mid\sigma^2=1,\ \deleted{\unsure{u}}\sigma u\sigma^{-1}=u^{-1}\}
\]
i.e.
\[
G_1=\underbrace{\mathbb Z.\mathbb Z/2\mathbb Z}_{\substack{\Vert\ \mathrm{dfn}\\D_\infty\\\text{(groupe diédral infini)}}}
\]
\margin{produit $\frac12$ direct, avec $\sigma$ (gén. de $\mathbb Z/2\mathbb Z$) opérant sur $\mathbb Z$ par symétrie}

La condition que \inserted{les involutions} $\sigma_0$ et $\sigma_1$ \deleted{\unsure{op}} opèrent ss pt
fixe dans $R$ s'écrivent par
\[
\left\{\begin{aligned}\sigma x&\ne x\\u^{-1}x&\ne\sigma x\end{aligned}\right.\qquad\forall x\in R
\]
La condition $\sigma_0x\ne\sigma_1x\quad\forall x\in R$ se transcrit par
\[
\bigm|\,u^{-1}x\ne x\qquad\forall x\in R
\]
donc ensemble avec les deux précédentes
cela signifie :
\[
\forall x\in R,\qquad\mathrm{Card}\{x,\sigma x,ux\}=3
\]
Noter que l'on a
\diagram{\begin{tikzcd}[ampersand replacement=\&,column sep=small]
1\arrow[r]\& \underbrace{G_1^{\unsure{0|\circ}}}_{\unsure{\simeq\mathbb Z|\mathbb Z}}\arrow[r]\& G_1\arrow[r]\&\deleted{\unsure{G}}\mathbb Z/2\mathbb Z\arrow[r]\&1
\end{tikzcd}}
Noter que $G_1\to G_1^{0}$ est formé d'él. d'ordre 2, \uline{chacun} définit splitting
% archive sheet number: 55
\scanpage{57}
\textcircled{\small 4} \uline{Cor. de Prop} Les polygônes comb. (= contours
connexes) correspondent aux espaces homogènes
sous $G_1=D_\infty$, tels que le stabilisateur d'un
point quelc. ne contienne ni $\sigma$, ni $u$, ni $\sigma u$

Quels sont ces sous-groupes \inserted{de $D_\infty$} ? \deleted{\unsure{Il est facile de}} \inserted{Je dis que ce}
\inserted{sont exactement les} sous-groupes \deleted{\unsure{en question contiennent|en question contenus et|qui ne soient contenus}} soit réduits
à 1, soit de la forme $n\mathbb Z$ \inserted{($n\in\mathbb N$, $n\geq2$)} (sous-groupe cyclique
infini engendré par $u^n$) \deleted{\unsure{-- dans le g|-- dans le}} c'est
donc un ss-groupe \uline{invariant} tel que le
quotient soit isomorphe au groupe
\[
\begin{aligned}
D_n&=\{\sigma,u\mid\sigma^2=1,\sigma u\sigma^{-1}=u^{-1},\deleted{\unsure{u}}u^n=1\}\\
&=\{\sigma_0,\sigma_1\mid\sigma_0^2=\sigma_1^2=\deleted{\unsure{1,}}(\sigma_0\sigma_1)^n=1\}\\
&\simeq \mathbb Z_n.\mathbb Z/2\mathbb Z\quad(\text{produit }\tfrac12\text{ direct, }\ldots.)
\end{aligned}
\]
\[
\text{où}\qquad 2\leq n\leq+\infty\qquad(n\in\mathbb N).
\]
\inserted{Condition évidemment suffisante. Pour la néc.}
Il suffit de prouver que \qquad a) $H\subset\deleted{\unsure{\mathbb Z}}G_1^{\unsure{0|\circ}}=\mathbb Z$
b) $H\ne\deleted{\unsure{G}}\unsure{\mathbb Z_1^0|G_1^0|\mathfrak G_1^\circ}$ i.e. $H\not\ni u\deleted{\unsure{\sigma}}$ (ce qui est vrai par
hyp. $\sigma,u,\sigma u\notin gHg^{-1}$ \inserted{$\forall g\in G_1$}). Si on avait $H\not\subset G_1^0$,
i.e. $\exists n\in\mathbb N$ avec
\[
u^n\sigma\in H
\]
Or on a \deleted{$\unsure{u^m\sigma|u^2\sigma}$}
\[
\left\{\begin{aligned}
\mathrm{int}(u^m)(\sigma)&=u^m\sigma u^{-m}=u^m(u^m\sigma)=u^{2m}\sigma\\
\mathrm{int}(u^m)(u\sigma)&=u^{2m+1}\sigma
\end{aligned}\right.
\]
ce qui implique que, ou bien $u^n\sigma$ conjugué
à $\sigma$ ($n$ pair) ou bien $u^n\sigma$ conjugué à $u\sigma$
($n$ impair), donc $\sigma$ ou $u\sigma$ appartient à un
conjugué de $H$, contrairement à l'hyp.
% archive sheet number: 56
\scanpage{58}
Soit $\Pi_n$ ($n\in\mathbb N$, $2\leq n\leq+\infty$) le polygône
combinatoire défini par $\unsure{D_n|\mathbb D_n}$, considéré comme
espace homogène sous $G_1$. Alors
\[
\Pi_n=(S_n,A_n,R_n)
\]
avec
\[
\underset{\substack{\Vert\\\{s_i\}_{i\in\mathbb Z_n}}}{S_n}=\deleted{\unsure{D_n/\sigma_1|D_n}}\underset{(\overline{\sigma_0})}{u\sigma}\backslash D_n\simeq\mathbb Z_n=
\]
[car $(u\sigma{\unsure{]|)}}^2=1$ \unsure{ou|or|donc} $u\sigma$ définit une section
de l'ext. de $\mathbb Z_2$ par $\mathbb Z_n$]
\[
\underset{\substack{\Vert\\\{a_i\}_{i\in\mathbb Z_n}}}{A_n}=\underset{\substack{\Vert\\(\sigma_1)}}{\sigma}\backslash D_n\simeq\mathbb Z_n
\]
Enfin, la relation d'incidence est donnée
par \qquad $(s_i,a_i)$ et $(s_{i+1},a_i)$ incidents.
\figure{\begin{tikzpicture}[x=1cm,y=1cm,baseline]
\draw (0,0) -- (1.55,.09) -- (2.95,1.04) -- (3.03,2.02);
\fill (1.55,.09) circle (.04) (2.95,1.04) circle (.04);
\node[below left,inner sep=1pt] at (0,0) {$s_0$};
\node[above,inner sep=1pt] at (.69,.09) {$a_0$};
\node[below right,inner sep=1pt] at (1.55,.09) {$s_1$};
\node[above left,inner sep=2pt] at (2.26,.65) {$a_1$};
\node[right,inner sep=2pt] at (2.95,1.04) {$s_2$};
\end{tikzpicture}}
On retrouve la description des polyg. comb.
\unsure{standart|standard|d'avant}, et le fait que la classe d'iso est
caractérisée par le nb des côtés = nb d'arêtes

\uline{\unsure{(Re)corollaire|Corollaire} 2} Si $\Pi,\Pi'$ sont deux polygônes
combinatoires iso, \deleted{alors} et $r\in\mathrm{Rep}(\Pi)$,
alors
\diagram{\begin{tikzcd}[ampersand replacement=\&,row sep=small,column sep=large]
\mathrm{Isom}(\Pi,\Pi')\arrow[r,"\sim"]\&\mathrm{Rep}(\Pi')\\
\varphi\arrow[r,mapsto]\&\varphi(r)
\end{tikzcd}}
C'est un cas particulier d'un \unsure{énoncé|théorème} sur les
% archive sheet number: 57
\scanpage{59}
espaces homogènes réguliers (ou ``distingués'').
\uline{Egalement} \unsure{?} :

\uline{Corollaire 3} Un polygône combinatoire de \unsure{rang $n$|rang|$n$-gone}
($2\leq n\leq\infty$) \deleted{muni} \uline{\unsure{de'un|de un|d'un} repère} est défini à
isom unique près -- et est \uline{can} \uline{isom}
à $\Pi_n$.

\uline{Corollaire 4}
\[
\begin{aligned}
\mathrm{Aut}(\Pi_n)&\simeq\unsure{D_n|\mathbb D_n}\quad(\text{can\textsuperscript{t}})\\
\mathrm{Aut}(\Pi)&\simeq D_n\quad(\text{non can\textsuperscript{t}})\ \text{si }\Pi\ n\text{-gône}
\end{aligned}
\]

\uline{NB} La restriction supplémentaire $\sigma_0x\ne\sigma_1x\quad\forall x\in E$
(en plus de $\sigma_0,\sigma_1$ ss pt fixe) sur les
\deleted{\unsure{espaces|orbites}} \inserted{$G_1$-ens} correspondants à des contours,
revient à l'\deleted{\unsure{élimination}} \inserted{exclusion} des orbites de la
forme $\deleted{\unsure{\Pi_2\simeq|\Pi_2=|\mathbb Z_2}}D_1$ -- \unsure{géom.|géos.} exclusion des
polygônes à 1 côté. On pourrait les
inclure, en n'exigeant pas $R\hookrightarrow S\times A$.

De ttes façons, = avec la restriction $n\geq 2$
pour les orbites, il semble qu'on soit encore
trop large dans la définition, si on compare
à celle des polyèdres \uline{combinatoires} -- où
on a exigé que \inserted{dans} l'ens ordonné associé
\deleted{deux} él. majorés aient un Sup. Ici \deleted{\unsure{ceti|cette|crois}}
l'exigence analogue reviendrait à exiger que
pour tte orbite on ait $n\geq3$ -- c'est à cette
% archive sheet number: 58
\scanpage{60}
condition, rappelons-le, que l'interprétation
d'un polygône comme une structure
\uline{sur $S$} (ens des sommets) -- savoir
un ens \unsure{de des|de de|des} 2 \unsure{paires|pourv.} \unsure{circulaire -- est|circulairement|circulaires -- est}
valable. \unsure{Par|Pour} la suite, on n'aura
à regarder que les cas $n\geq3$.

6. \uline{Orientation} d'un contour

a) Soit $R^+\subset R$, \deleted{\unsure{on a|on}} $R^-=R\smallsetminus R^+$. Conditions
équivalentes
\[
\begin{aligned}
\text{a)}\quad&\sigma_0R^+=\sigma_1R^+=R^-\\
\text{b)}\quad&R^+\xrightarrow{\sim} S,\quad R^+\simeq A
\end{aligned}
\quad\left|\begin{gathered}\text{On dit alors que }R^+\\\text{est une orientation}\end{gathered}\right.
\]
c'est ainsi une partie \inserted{$R^+$} de $R\deleted{\unsure{{}^+}}$ qui est le graphe
d'une \uline{bijection} $S\overset{\sim}{\longleftrightarrow} A$.
\margin{Orientation \deleted{\unsure{positive|opposée}} \inserted{opposée} d'une orientation \deleted{\unsure{déf}} (involution -- \unsure{ss|est} $\unsure{\ne\mathrm{id}|\ne 0}$ si $\Pi\ne\varnothing$)}
\[
\text{b)}\qquad\mathrm{Or}\left(\coprod_{i\in I}\deleted{\unsure{R_i(C)|R_i(C_i)}}\underset{\deleted{\text{contour}}}{C_i}\right)\simeq\prod_i\mathrm{Or}(C_i)
\]
si $(C_i)_{i\in I}$ est une famille de contours.

c) Si $\Pi$ est un polygône combinatoire,
alors $\mathrm{Card}\,\mathrm{Or}(\Pi)=2$, et $R$ est \inserted{alors} somme
disjointe de la \deleted{\unsure{ens|en}}famille des deux orientations
de $\Pi$, \deleted{\unsure{Pour tte orientation}} \inserted{lesquelles sont opposées l'une de l'autre.} Les deux \unsure{orient\textsuperscript{ns}|orientations}
sont les deux orbites de $\unsure{G_1^{0}|G_1^\circ|\mathfrak G_1^\circ}$ dans $R$.
% archive sheet number: 59
